%!TEX root = ./../../springer-mono-format/book.tex
%!TEX root = ../../thesis-trier.tex

Synchronizing automata were already defined in Section~\ref{intro:sec:basic}. Here, we collect further properties
that we need in later chapters.
\Cref{prelim:classes:sync:ideal} states that the set of synchronizing words is an ideal language.
\Cref{prelim:sync:trap} investigates synchronizing automata with a trap state.
\Cref{prelim:classes:sync:subfactor} relates synchronizing words to subautomata, morphic images,
quotient automata and vice versa.
\Cref{prelim:classes:sync:reach} 
investigates the set of maximal states in a synchronizing automaton.
\Cref{prelim:sync:pair} expresses synchronization in terms of pairs of states.
Lastly, \cref{prelim:sync:algorithmic} collects algorithmic results related
to practical implementations (results from a more theoretical perspective in the framework
of computational complexity theory are collected in \cref{chap:prelim_complexity_results} later).

\todo[inline]{genauer ergebnisse nennen? überblick?}

%\subsection{The Set of Synchronizing Words and Regular Ideal Languages} 
\subsection{Sets of Synchronizing Words and Ideal Languages} 
\label{prelim:classes:sync:ideal}

%!TEX root = ../../../springer-mono-format/book.tex
%!TEX root = ../../../thesis-trier.tex

The set of synchronizing words of a semiautomaton
$\semiautomat{A}{Q}{\Sigma}{\delta}$ is a regular ideal language:
it is an ideal, since $\delta(Q, uwv) \subseteq \delta(\delta(Q,w), v)$
is a singleton whenever $\delta(Q,w)$ is, and it is regular, since
$\Syn(\mathcal A) = L(\mathcal P(\mathcal A)_{Q, \mathcal F})$
with $\mathcal F = \{ \{q\} \mid q \in Q \}$.

Next, suppose $\automat{A}{Q}{\Sigma}{\delta}{q_0}{F}$ is minimal and accepts a
{non-empty} ideal language. Then, as observed by
Maslennikova~\cite{DBLP:journals/corr/Maslennikova14,DBLP:journals/ijfcs/Maslennikova19},
\begin{equation}\label{prelim:eqn:reg_ideal}
	L(\mathcal A) = \Syn(\mathcal A).
\end{equation}
The assumption $L(\mathcal A) \ne \emptyset$ cannot be dropped: the minimal
automaton for $\emptyset$ is $(\{q\}, \Sigma, \delta, q, \emptyset)$ with
$\delta(q,x) = q$, which is synchronizing with $\Syn(\mathcal A) = \Sigma^*$.
Of course, $\emptyset$ is still the set of synchronizing words of any
non-synchronizing semiautomaton, and so we obtain the following
characterization.

\begin{theorem}
	\label{prelim:thm:sync_words_is_ideal}
	Let $I \subseteq \Sigma^*$. Then $I$ is a regular ideal language if and only
	if $I = \Syn(\mathcal A)$ for some finite semiautomaton
	$\semiautomat{A}{Q}{\Sigma}{\delta}$.
\end{theorem}

By \cref{prelim:thm:sync_words_is_ideal}, investigating acceptors for sets of
synchronizing words, as we do in \cref{part:aut}, amounts to investigating
acceptors for regular ideal languages. In the following, we use the fact that
$\Syn(\mathcal A)$ is an ideal language without further mention.

As a shortest word in a regular language accepted by an $n$-state
automaton has length at most $n-1$, \cref{prelim:eqn:reg_ideal} yields the
following bound.

\begin{proposition}
	\label{prelim:prop:reset_length_ideal}
	If $\automat{A}{Q}{\Sigma}{\delta}{q_0}{F}$ is a minimal $n$-state automaton
	accepting an ideal language, then $\mathcal A$ is synchronizing with reset
	length at most $n-1$.
\end{proposition}

For $i \ge 1$, the minimal automaton for $\Sigma^* a^i \Sigma^*$ has $i+1$ states
and reset threshold $i$, so the bound is sharp for every $n \ge 1$.

\subsection{Subautomata, Morphic Images, and Quotient Automata}
\label{prelim:classes:sync:subfactor}

%!TEX root = ../../../thesis-trier.tex

\Cref{prelim:prop:sync_word_factor} collects statements on
how the set of synchronizing words related to subautomata, surjective morphic images,
and quotient automata. \Crefitem{prelim:prop:sync_word_factor}{prelim:Synuv}
is due to Rystsov~\cite[Prop.~2.3]{Rystsov97}.

\begin{proposition}
	\label{prelim:cor:fact_sync_word}
	\label{prelim:prop:morphism_sync_word}
%	\label{constr_input:prop:rho_trap_sync_factor}
	\label{prelim:prop:sync_word_factor}
	\label{prelim:prop:sync_words_restriction}
	Let $\semiautomat{A}{Q}{\Sigma}{\delta}$. Then the following is true.
	
	\begin{enumerate}
		\item \label{prelim:Synsurj} $\Syn(\mathcal A) \subseteq \Syn(\mathcal B)$ if we have a surjective morphism $\varphi \colon Q \to S$ onto $\semiautomat{B}{S}{\Sigma}{\mu}$.
		
		\item \label{prelim:Syncongr} $\Syn(\mathcal A) \subseteq \Syn(\mathcal A / {\sim})$ for a $\TS{\mathcal A}$-congruence $\sim$ on $\mathcal A$.
		
		\item \label{prelim:Synsubaut} $\Syn(\mathcal A) \subseteq \Syn(\mathcal A_{|S})$ for an upper set $S$.  So, 
		if $\mathcal A$ is synchronizing, then $\mathcal A_{|S}$ is synchronizing.
		
		\item \label{prelim:Synuv} $\Syn(\mathcal A / {\sim_S}) \Syn(\mathcal A_{|S}) \subseteq \Syn(\mathcal A)$ for an upper set $S$ (Rystsov~\cite[Prop.~2.3]{Rystsov97}).
%		More specifically, 
%		for $u \in \Syn(\mathcal A / {\sim_S})$ and $v \in \Syn(\mathcal A_{|S})$ 
%	    it is $\applytrans{\delta}{Q}{uv} = \applytrans{\delta}{S}{v}$.
		
	%	\item \label{prelim:Synrho} $\Syn(\mathcal A / \eqreach) \subseteq \Syn(\mathcal A)$ for a $\rho$-automaton $\mathcal A$ with a trap state.
	\end{enumerate}
\end{proposition}
\begin{proof}
	\ref{prelim:Synsurj} Suppose $|\delta(Q, u)| = 1$ for $u \in \Sigma^*$.
	We have $\mu(S, u) = \mu(\varphi(Q), u) = \varphi(\delta(Q, u))$.
	So $|\mu(S, u)| = 1$.	
	
	\ref{prelim:Syncongr} Applying~\ref{prelim:Synsurj} 
	to the canonical morphism for $\sim$
	from  \cref{prelim:prop:congruence_induces_morphism}
	gives the claim.
	
%	By \cref{prelim:prop:congruence_induces_morphism}
%	the canonical projection for $\sim$ is morphism.
%	Then~\ref{prelim:Synsurj} 
%    gives the claim.	
	
	
%	This follows from~\ref{prelim:Synsurj}
%	applied to the canonical morphism for $\sim$
%	and \cref{prelim:prop:congruence_induces_morphism}.
	
	\ref{prelim:Synsubaut} The equation $\delta_{|S}(S, u) \subseteq \delta(Q, u)$ for $u \in \Sigma^*$
	gives the claim.
	
	\ref{prelim:Synuv} Let $u \in \Syn(\mathcal A/{\sim_S})$.
	The set $S$ is a trap state in $\mathcal A / {\sim_S}$.
	So, this state must be the synchronizing state in $\mathcal A / {\sim_S}$, i.e., we have $\applytrans{\delta_{\sim_S}}{[q]_{\sim_S}}{u} = S$ for all $q \in Q$, or equivalently $\applytrans{\delta}{Q}{u} \subseteq S$.
	Then for $v \in \Syn(\mathcal A_{|S})$ we get 
	$\applytrans{\delta}{Q}{uv} = \applytrans{\delta}{\applytrans{\delta}{Q}{u}}{v} \subseteq \applytrans{\delta}{S}{v}$ where $|\delta(S, v)| = 1$.
	Hence $uv \in \Syn(\mathcal A)$.
\end{proof}

\begin{corollary} \label{prelim:syncaut:cor:Syncongr}
	Let $\semiautomat{A}{Q}{\Sigma}{\delta}$ 
	be an automaton with a trap state and a $\TS{\mathcal A}$-congruence~$\sim$ on its states
	such that  $[t]_{\sim} = \{t\}$ for each trap state~$t$.
	Then we have $\Syn(\mathcal A) = \Syn(\mathcal A / {\sim})$.
\end{corollary} 
\begin{proof}
	By \crefitem{prelim:cor:fact_sync_word}{prelim:Syncongr} 
	we only have to prove that $ \Syn(\mathcal A / {\sim})$
	is contained in $\Syn(\mathcal A)$.
	So, let $u \in \Syn(\mathcal A / {\sim})$.
	Let~$t$ be some trap state in~$\mathcal A$.
	By \cref{prelim:aut:prop:morph:trap-to-trap},
	also $[t]_{\sim}$ is a trap state in $\mathcal A / {\sim}$.
	By \cref{prelim:thm:sync_trap_state},
	we have $\delta_{\sim}([q]_{\sim}, u) = [t]_{\sim}$
	for all $q \in Q$,	i.e., $\delta(q, u) \sim t$ for all $q \in Q$.
	Since $[t]_{\sim} = \{t\}$
	this yields $\delta(q, u) = t$ for all $q \in Q$
	and so $u \in \Syn(\mathcal A)$.	
\end{proof}

\begin{corollary} \label{lem:sub-quotient-extremal-linear}
	Let $\semiautomat{A}{Q}{\Sigma}{\delta}$
	be a synchronizing $n$-state automaton
	all whose synchronizing words have length at least $n - 1$
	and $S \subseteq Q$ be an upper set. 	
	Then all synchronizing words 
	of $\mathcal A_{|S}$ and $\mathcal A / {\sim_S}$
	have length at least $|S| - 1$
	and $n - |S|$, respectively.
\end{corollary}
\begin{proof}
	If $\mathcal A_{|S}$
	or $\mathcal A / {\sim_S}$
	have a synchronizing word
	of shorter length,
	then by \crefitem{prelim:cor:fact_sync_word}{prelim:Synuv}
	also $\mathcal A$ would have a synchronizing
	word of shorter length.	
\end{proof}

\subsection{Synchronizing Automata with a Trap State} 
\label{prelim:sync:trap}

%!TEX root = ../../../springer-mono-format/book.tex
%!TEX root = ../../../thesis-trier.tex

Several classes of automata we investigate later have the property that a
synchronizing state must be a trap state. Here, we collect statements about
synchronizing automata with a trap state.
First, \cref{prelim:thm:sync_trap_state} characterizes synchronizability for
automata with a trap state. It also yields a sharp bound on the reset length of
synchronizing automata with a trap state, first stated by Rystsov~\cite{Rys77}
(see also~\cite[Thm.~6.1]{Rystsov97}) and
Babcs{\'{a}}nyi~\cite[Cor.~3]{DBLP:conf/birthday/Babcsanyi09}; we state this
bound, together with a family of automata attaining it, in
\cref{prelim:rem:trap_reset_length}.
Then, in \cref{prelim:cor:sync_trap_state}, we state some consequences. For
example, we characterize the synchronizing automata with a trap state as those
automata possessing a unique maximal state with respect to the reachability order.


\begin{theorem}
	\label{prelim:thm:sync_trap_state}
	Let $\semiautomat{A}{Q}{\Sigma}{\delta}$ be an $n$-state automaton with a trap state $t \in Q$.
	Then the following are equivalent.
	
	\begin{enumerate}
		\item \label{prelim:subclasses:thm:trap:item:sync} $\mathcal A$ is synchronizing.
		
		\item \label{prelim:subclasses:thm:trap:item:syncunique} $\mathcal A$ is synchronizing with unique synchronizing state $t$.
		
		\item \label{prelim:subclasses:thm:trap:item:unique-max} $\forall q \in Q \  \exists u \in \Sigma^* \colon \applytrans{\delta}{q}{u} = t$, i.e., $q \lereach t$ for all $q \in Q$.
		
		\item \label{prelim:subclasses:thm:trap:item:bound} $\exists u \in \Sigma^* \colon |u| \le \frac{n(n-1)}{2} \land ( \forall q \in Q \colon \applytrans{\delta}{q}{u} = t )$. 
	\end{enumerate}	
\end{theorem}
\begin{proof}
	$\ref{prelim:subclasses:thm:trap:item:sync} \Rightarrow
\ref{prelim:subclasses:thm:trap:item:syncunique}$:
Let $\applytrans{\delta}{Q}{u} = \{q\}$ for some $u \in \Sigma^*$ and
$q \in Q$. As $t$ is a trap state, $t \in \applytrans{\delta}{Q}{u}$, so
$t = q$. Hence $\{t\}$ is the set of synchronizing states.
	
	$\ref{prelim:subclasses:thm:trap:item:syncunique} \Rightarrow \ref{prelim:subclasses:thm:trap:item:unique-max}$: By assumption there exists $u \in \Sigma^*$ 
	such that $\applytrans{\delta}{q}{u} = t$ for all $q \in Q$, i.e., $q \lereach t$
	for all $q \in Q$.	
	
	$\ref{prelim:subclasses:thm:trap:item:unique-max} \Rightarrow \ref{prelim:subclasses:thm:trap:item:bound}$: 	
	For $n = 1$ we have $Q = \{t\}$ and $\varepsilon$ is a synchronizing word.
	Thus, assume $n > 1$ in the following.
	
	We construct words $u_i$ and sets $Q = Q_0, Q_1, \ldots$
	containing~$t$ with strictly decreasing cardinalities as follows. Set $Q_0 = Q$.
	Given $Q_i$ with $t \in Q_i$, stop if $Q_i = \{t\}$.
	Otherwise, since $t$ is reachable from each state by \cref{prelim:subclasses:thm:trap:item:unique-max},
	we can pick a shortest word $u_{i+1}$ with $\delta(q_i, u_{i+1}) = t$
	for some $q_i \in Q_i \setminus \{t\}$. Then set $Q_{i+1} = \delta(Q_i, u_{i+1})$.
	
	Since~$t$ is a trap state, we have $\delta(t, u_{i+1}) = t = \delta(q_i, u_{i+1})$.
	So $t \in Q_{i+1}$ and $|Q_{i+1}| < |Q_i|$.
	Thus there exists $m \le n - 1$ such that $Q_m = \{t\}$.
	
	Let $i \in \{0,1,\ldots,m-1\}$ and 
	let $q_i = p_0, p_1, \ldots, p_\ell = t$ be the path induced by $u_{i+1}$,
	where $\ell = |u_{i+1}|$. By the minimal choice of~$u_{i+1}$, no $p_r$ with
	$r \ge 1$ lies in $Q_i \setminus \{t\}$ and no $p_r$ with $r < \ell$
	equals~$t$, and the $p_r$ are pairwise distinct; otherwise a proper suffix,
	a proper prefix, or the word obtained by deleting a factor of~$u_{i+1}$
	would take a state of $Q_i \setminus \{t\}$ to~$t$. As
	$p_1, \ldots, p_{\ell-1}$ are thus $\ell - 1$ distinct states outside~$Q_i$,
	of which there are exactly $n - |Q_i|$, we get
	$|u_{i+1}| \le n - |Q_i| + 1$.
	
	Set $u = u_1 \cdots u_m$. We have $\delta(Q, u) = \{t\}$.
	As $n = |Q_0| > |Q_1| > \cdots > |Q_{m-1}| \ge 2$, the numbers
	$n - |Q_{i-1}| + 1$ for $i \in \{1, \ldots, m\}$ are pairwise distinct and
	lie in $\{1, \ldots, n-1\}$.
	Thus, we have	
	\[
	|u| = |u_1| + \cdots + |u_m| \le \sum_{i=1}^m \bigl( n - |Q_{i-1}| + 1 \bigr)
	\le \sum_{i=1}^{n-1} i = \frac{n(n-1)}{2}.
	\]
	
		
	$\ref{prelim:subclasses:thm:trap:item:bound} \Rightarrow \ref{prelim:subclasses:thm:trap:item:sync}$: This is obvious.
\end{proof}

\begin{remark} \label{prelim:rem:trap_reset_length}
	By \cref{prelim:thm:sync_trap_state}, every synchronizing $n$-state automaton with a trap state has reset length
	at most $\frac{n(n-1)}{2}$, and this bound is sharp: the automata
	from \cref{prelim:fig:trap_state_tight} have reset length
	exactly $\frac{n(n-1)}{2}$.
\end{remark}

\begin{figure}[htb]
	\centering
	\begin{tikzpicture}[node distance = 4em and 4em]
		\node[state]                  (0)     {$1$};
		\node[state, right = of 0]    (1)     {$2$};
		\node[state, right = of 1]    (2)     {$3$};
		\node[right = of 2]           (dots)  {$\cdots$};
		\node[state, right = of dots] (n2)    {$n$};
		
		\path[->] (1) edge [above] node {$a_1$} (0)
		(1)    edge [bend left, above] node {$a_2$}     (2)
		(2)    edge [bend left, below] node {$a_2$}     (1)
		(2)    edge [bend left, above] node {$a_3$}     (dots)
		(dots) edge [bend left, below] node {$a_3$}     (2)
		(dots) edge [bend left, above] node {$a_{n-1}$} (n2)
		(n2)   edge [bend left, below] node {$a_{n-1}$} (dots);
	\end{tikzpicture}
	\caption[Automata with a trap state attaining the maximal reset length]{Trap-state automata
		from~\cite{Rystsov97} with reset
		length $\frac{n(n-1)}{2}$. The letter $a_1$ maps $2$ to the trap state~$1$,
		and $a_i$ interchanges $i$ and $i+1$ for $i \ge 2$; self-loops are omitted.
		A shortest synchronizing word is	
		$u = w_1 w_2 \cdots w_{n-1}$
		where $w_i = a_i a_{i-1} \cdots a_1$ for $i \in \{1, \ldots, n-1\}$.		
		Note that $\Sigma = \{a_1, a_2, \ldots, a_{n-1}\}$ grows with~$n$.}
	\label{prelim:fig:trap_state_tight}
\end{figure}

\Cref{prelim:subclasses:thm:trap:item:syncunique} of
\cref{prelim:thm:sync_trap_state} states that in a synchronizing automaton, a
trap state is the unique synchronizing state.

\begin{corollary}
	\label{prelim:classes:sync:cor:sync-implies-at-most-one-trap}
	In a synchronizing automaton, there exists at most one trap state.
\end{corollary}
\begin{proof}
	Let $t, t'$ be trap states. By \cref{prelim:subclasses:thm:trap:item:syncunique}
	of \cref{prelim:thm:sync_trap_state},
	each of them is the unique synchronizing state, so $t = t'$.
\end{proof}

Thus, by \cref{prelim:thm:sync_trap_state}, an automaton with a trap state is
synchronizing if and only if that trap state is reachable from every state. The
trap state is then unique by
\cref{prelim:classes:sync:cor:sync-implies-at-most-one-trap}. Note, however,
that it is easy to find automata with a unique trap state that are not
synchronizing.


\begin{corollary} \label{prelim:cor:sync_trap_state}
	\label{prelim:classes:sync:cor:sync-trap-state}
Let $\semiautomat{A}{Q}{\Sigma}{\delta}$. Then the following are equivalent.
\begin{enumerate}
\item \label{prelim:cor:sync_trap_state:item:sync} $\mathcal A$ is synchronizing with a trap state.
	
\item \label{prelim:cor:sync_trap_state:item:trapreach} There exists a trap state $t \in Q$ such that $q \lereach t$ for all $q \in Q$.

\item \label{prelim:cor:sync_trap_state:item:uniquemaxclass} $Q / \eqreach$ has a maximum, and this maximum is a singleton $\{t\}$ for some $t \in Q$.

\item  \label{prelim:cor:sync_trap_state:item:uniquemax} $\mathcal A$ has a unique maximal state.

\item \label{prelim:cor:sync_trap_state:item:uniquesync} $\mathcal A$ has a unique synchronizing state.
\end{enumerate}
\end{corollary}
\begin{proof}
$\ref{prelim:cor:sync_trap_state:item:sync}\Rightarrow \ref{prelim:cor:sync_trap_state:item:trapreach}$: This follows from \cref{prelim:subclasses:thm:trap:item:unique-max} of \cref{prelim:thm:sync_trap_state}.

$\ref{prelim:cor:sync_trap_state:item:trapreach}\Rightarrow \ref{prelim:cor:sync_trap_state:item:uniquemaxclass}$: Since $t$ is a maximum, $[t]_{\eqreach}$ is a maximum in $Q / {\eqreach}$.
 Since~$t$ is a trap state, we have $[t]_{\eqreach} = \{t\}$.

$\ref{prelim:cor:sync_trap_state:item:uniquemaxclass}\Rightarrow \ref{prelim:cor:sync_trap_state:item:uniquemax}$: Since~$\{t\}$
is a maximum in $Q / {\eqreach}$, $t$ is maximal.
Let $q \in Q$ be a maximal state. By assumption, $q \lereach t$.
 By maximality, $t \lereach q$. Thus $q \eqreach t$, which gives $q = t$ by assumption.
 
$\ref{prelim:cor:sync_trap_state:item:uniquemax}\Rightarrow \ref{prelim:cor:sync_trap_state:item:uniquesync}$: Let $t$ be the unique maximal state. By \cref{prelim:order:maximal-elements-upper-set},
 the set of maximal states $\{t\}$ is an upper set. Thus, $t$ is a trap state.
 Since every state is below some maximal state by finiteness, and there is only one maximal
 state, we have $q \lereach t$ for each $q \in Q$.
 Thus, by \cref{prelim:thm:sync_trap_state}, $t$ is a unique synchronizing state.
 
$\ref{prelim:cor:sync_trap_state:item:uniquesync}\Rightarrow \ref{prelim:cor:sync_trap_state:item:sync}$: Let $t$ be the unique synchronizing state. Then $\delta(Q, u) = \{t\}$
for some $u \in \Sigma^*$. Let $a \in \Sigma$. We have $|\delta(Q, ua)| = 1$.
By assumption, $\delta(Q, ua) = \{t\}$.
Thus $\delta(t, a) \in \delta(Q, ua) = \{t\}$, and so $t$ is a trap state. 
\end{proof}

\subsection{Synchronizing Automata and the Reachability Order}
\label{prelim:classes:sync:reach}

%!TEX root = ../../../thesis-trier.tex

Here, %we investigate synchronizing automata in terms of the reachability order on the states
%and the set of maximal states for this quasiorder.
in \cref{prelim:prop:maximal-states-factor-sync},
we first prove that the quotient automaton w.r.t.~the set of maximal
states is synchronizing.
Then, in \cref{prelim:thm:SCC_sync_words}, we prove that an automaton is synchronizing if and only if the subautomaton
on the set of maximal states is synchronizing and strongly connected.
Further, in this case, the set of maximal states equals the set of synchronizing states.
As an application, in \cref{prelim:prop:CernySCC},
we show that if the \v{C}ern\'y conjecture holds for strongly connected automata,
then it holds for all automata.

\begin{proposition}
	\label{prelim:prop:maximal-states-factor-sync}
	Let $\semiautomat{A}{Q}{\Sigma}{\delta}$ and $S \subseteq Q$ be
	the set of maximal states for $\lereach$.
	Then~$S$ is an upper set and $\mathcal A / {\sim_S}$ is synchronizing with trap state~$S$.
\end{proposition}
\begin{proof}
   By \cref{prelim:order:maximal-elements-upper-set}, $S$ is an upper set forming a trap state in $\mathcal A / {\sim_S}$. 
   By \cref{prelim:automata:prop:max-states}, this trap state is reachable from every other 
   state, and so the claim follows with
   \crefitem{prelim:cor:sync_trap_state}{prelim:cor:sync_trap_state:item:trapreach}.
\end{proof}

\begin{theorem}
	\label{prelim:thm:SCC_sync_words}
	Let $\semiautomat{A}{Q}{\Sigma}{\delta}$ and $S \subseteq Q$ be the set of maximal states for $\lereach$.
	Then $S$ is an upper set and the following are equivalent.
	
	\begin{enumerate}
		\item \label{prelim:thm:SCC_sync_words:sync} $\mathcal A$ is synchronizing.
		
		\item \label{prelim:thm:SCC_sync_words:Ssync} $\mathcal A_{|S}$ is synchronizing.
		
		\item \label{prelim:thm:SCC_sync_words:SsyncSCC} $\mathcal A_{|S}$ is synchronizing and strongly connected (so that $S$ is a sink component).
		
		\item \label{prelim:thm:SCC_sync_words:Ssyncstates} $\mathcal A$ is synchronizing and $S$ equals the set of synchronizing states of $\mathcal A$.
	\end{enumerate}
\end{theorem}
\begin{proof} 
	By \cref{prelim:order:maximal-elements-upper-set}, $S$ is an upper set. Next, we show the stated equivalences.	
	
	$\ref{prelim:thm:SCC_sync_words:sync} \Rightarrow~\ref{prelim:thm:SCC_sync_words:Ssync}$: By \crefitem{prelim:prop:sync_words_restriction}{prelim:Synsubaut}, 
	if $\mathcal A$ is synchronizing, then $\mathcal A_{|S}$
	is synchronizing.	
		
	$\ref{prelim:thm:SCC_sync_words:Ssync} \Rightarrow~\ref{prelim:thm:SCC_sync_words:SsyncSCC}$: 
	By assumption $\delta(S, u) = \{q\}$ for some $q \in Q$ and $u \in \Sigma^*$.
	So $p \lereach q$ for all $p \in S$,
	which gives, as $p$ is a maximal state, $p \eqreach q$ for all $p \in S$.
	Hence $S$ is strongly connected and as it is an upper set
	it is a sink component.
	
	$\ref{prelim:thm:SCC_sync_words:SsyncSCC} \Rightarrow~\ref{prelim:thm:SCC_sync_words:Ssyncstates}$: 
	By \cref{prelim:prop:maximal-states-factor-sync}
	the automaton $\mathcal A / {\sim_S}$ is synchronizing.
	So, from \crefitem{prelim:prop:sync_word_factor}{prelim:Synuv}
	we get that $\mathcal A$ is synchronizing.
	Let $q \in Q$ be such that $\delta(Q, u) = \{q\}$ for some $u \in \Sigma^*$.
	Then $p \lereach q$ for all $p \in Q$. So, $q$ is maximal and we have $q \in S$.
	Furthermore, as~$S$ is a sink component, for $p \in S$
	there exists $v \in \Sigma^*$ such that $\delta(q, v) = p$.
	Then $\delta(Q, uv) = \{p\}$ so that $p$ is a synchronizing state.
	So we obtain that $S$ equals the set of synchronizing states.			
	
	$\ref{prelim:thm:SCC_sync_words:Ssyncstates} \Rightarrow~\ref{prelim:thm:SCC_sync_words:sync}$: This is obvious.
\end{proof}

\begin{corollary} \label{prelim:classes:sync:cor:sync-state}
	Let $\semiautomat{A}{Q}{\Sigma}{\delta}$ and $q \in Q$.
	Then the following are equivalent.
	\begin{enumerate}
	\item $q$ is synchronizing.
	
	\item $\mathcal A$ is synchronizing and~$q$ is maximal w.r.t.~$\lereach$.
	
	\item $\mathcal A$ is synchronizing and~$q$ is a maximum w.r.t.~$\lereach$.
		
	\end{enumerate}
\end{corollary}

\crefitem{prelim:thm:SCC_sync_words}{prelim:thm:SCC_sync_words:SsyncSCC}
and \cref{prelim:prop:maximal-states-SCC} give the next corollary.

\begin{corollary}\label{prelim:cor:sync-unique-sink}
	A synchronizing automaton has a unique sink component.
\end{corollary} 

Furthermore, \cref{prelim:thm:SCC_sync_words}
yields the following corollary (but this also follows from \cref{prelim:autsyncw:lem:minrank} stated earlier).
%we obtain that strongly connected synchronizing automata
%are precisely those for which every state is synchronizing.
 
\begin{corollary}
	\label{prelim:cor:SCC_sync}
	In an automaton every state is a synchronizing state if and only if it is synchronizing and strongly connected.
\end{corollary} 

As a nice application of \cref{prelim:prop:sync_word_factor} and \cref{prelim:thm:SCC_sync_words}, 
we prove that confirmation of
the \v{C}ern\'y conjecture for strongly connected automata
implies it for all automata.
This result appears in the literature in~\cite[Thm.~10]{Pin81a}, \cite[Prop.~2]{DBLP:journals/acta/CarpiD09}, \cite[Thm.~2]{Rystsov92}, and in slightly more general form in
\cite[Prop.~2.1]{DBLP:journals/tcs/Volkov09} and~\cite[Thm.~6]{DBLP:conf/birthday/Babcsanyi09}.

\begin{proposition} \label{prelim:prop:CernySCC}
	The \v{C}ern\'y conjecture is true for all automata if  it is true for all strongly connected automata.
\end{proposition}
\begin{proof}
	Assume the \v{C}ern\'y conjecture is true for strongly connected automata.
	Let $\semiautomat{A}{Q}{\Sigma}{\delta}$ be a synchronizing $n$-state automaton
	with set of synchronizing states $S \subseteq Q$.
	By \cref{prelim:thm:SCC_sync_words} the set~$S$ equals the upper set of all maximal states
	and $\mathcal A_{|S}$ is synchronizing and strongly connected.
	So, by assumption,~
	$\mathcal A_{|S}$ has a synchronizing word $v \in \Sigma^*$ of length at most $(|S| - 1)^2$.
	%As $\mathcal A$ is not strongly connected, we have $|S| < n$.
	By \cref{prelim:prop:maximal-states-factor-sync}, $\mathcal A / {\sim_S}$ is synchronizing with trap state $S$.
	As~$\mathcal A / {\sim_S}$ has $n - |S| + 1$ states,
	from \crefitem{prelim:thm:sync_trap_state}{prelim:subclasses:thm:trap:item:bound}
    we obtain that $\mathcal A / {\sim_S}$ 
	has a synchronizing word $u \in \Sigma^*$ of length at most $\frac{( n - |S| + 1 )(n - |S|)}{2}$.	
	By \crefitem{prelim:prop:sync_word_factor}{prelim:Synuv}, the word~$uv$
	synchronizes~$\mathcal A$ and
	\begin{align*}
		|uv| \le \frac{( n - |S| + 1 )(n - |S|)}{2} + (|S| - 1)^2
		 & \le (n - |S|)^2 + (|S|-1)^2 \\
		 & \le (n - |S| + |S| - 1)^2 = (n - 1)^2. \qedhere
	\end{align*}
\end{proof}

Volkov~\cite[Prop.~2.1]{DBLP:journals/tcs/Volkov09} (see also~\cite[Prop.~1]{AnaVolGus2013})
gave the following generalization
of the previous argument. Suppose we have a class of automata
fulfilling the following two properties:
\begin{enumerate} 
	\item it is closed under taking
	subautomata and Rees quotient automata, and
	\item the reset length of every  $n$-state synchronizing automaton from this class that is strongly connected
	or has a trap state is bounded by~$f(n)$ for some function $f \colon \mathbb N_0 \to \mathbb N_0$
	with $f(n) + f(m) \le f(n + m - 1)$ for all $n,m \ge 1$.
	
%	there exists  a function 
%	$f \colon \mathbb N_0 \to \mathbb N_0$ with
%	$f(n) + f(m) \le f(n + m - 1)$
%	for all $n, m \ge 1$
%	such that the reset length of every $n$-state
%	automaton from this class that is strongly connected
%	or has a trap state is bounded by~$f(n)$.
\end{enumerate} 
Then, with the notation as in the proof of \cref{prelim:prop:CernySCC},
$|uv| \le f(n - |S| + 1) + f(|S|) \le f(n)$, and we obtain that
every $n$-state synchronizing automaton in this class
has a synchronizing word of length at most $f(n)$.
Specifically, for the class of all automata, the function $f(n) = (n-1)^2$ fulfills both conditions.

\subsection{The Pair Criterion} % Condition}
\label{prelim:sync:pair}

%!TEX root = ../../../thesis-trier.tex

The pair criterion (\cref{prelim:thm:pair_sync_condition}), due to \v{C}ern\'y~\cite[Thm.~2]{Cer64},
states that an automaton is synchronizing if and only if all pairs of states can be mapped to a single state.
Here, we deduce this important criterion from \THM~\ref{prelim:thm:sync_trap_state} and the observation, stated next,
that the pair automaton from \cref{sec:pair_aut}
is an automaton with a trap state having the same set of synchronizing states.

\begin{proposition}[{Rystsov~\cite[Prop.~5.2]{Rystsov97}}]
	\label{prelim:subclasses:sync:Aiffpairaut}
	Let $\semiautomat{A}{Q}{\Sigma}{\delta}$.
	Then we have $\Syn(\mathcal A) = \Syn\left(\mathcal A^{\le 2} / {\sim_F}\right)$ where $F = \{\{q\} \mid q \in Q\}$.
	So, $\mathcal A$ is synchronizing if and only if its pair automaton 
	$\mathcal A^{\le 2} / {\sim_F}$
	is synchronizing.
\end{proposition}
\begin{proof}
	Suppose $u \in \Syn(\mathcal A)$. Then $\applytrans{\delta}{Q}{u} = \{q\}$ for some $q \in Q$.
	Hence $\applytrans{\delta}{\{p_1, p_2\}}{u} = \{q\}$
    for all $p_1, p_2 \in Q$.
	So, $u$ maps all states of $\mathcal A^{\le 2}$
	into $F$, giving $u \in \Syn\left(\mathcal A^{\le 2} / {\sim_F}\right)$.	
	
	Conversely, let $u \in \Syn\left(\mathcal A^{\le 2} / {\sim_F}\right)$.
	Pick some $q \in Q$.
	Consider a state $p \in Q$.
	Since~$F$ is a trap state, it must be the unique synchronizing state
	and so $\delta_{\le 2}(\{p,q\}, u) \in F$.
	Thus~$\delta(p, u) = \delta(q, u)$
	and since~$p$ was chosen arbitrarily,
	we find $\delta(Q, u) = \{\delta(q, u)\}$
	and so~$u$ synchronizes~$\mathcal A$.
\end{proof}

We remark that it also easily follows that $\Syn(\mathcal A^{\le 2}) = \Syn(\mathcal A)$.

\begin{remark}
 Since the pair automaton has $\binom{n}{2} + 1$ states,
 by \cref{prelim:subclasses:sync:Aiffpairaut},
 the length of a shortest synchronizing words is quadratically bounded
 in the number of states if and only if 
 it is linearly bounded in the number of states of the pair automaton.	
\end{remark}

The pair automaton has $F$ as a trap state. By \cref{prelim:subclasses:sync:Aiffpairaut} and 
by applying the quantifier-swapping for trap-state synchronizing automata stated
in \cref{prelim:thm:sync_trap_state} \ref{prelim:subclasses:thm:trap:item:unique-max} and~\ref{prelim:subclasses:thm:trap:item:bound} 
to $\mathcal A^{\le 2} / {\sim_F}$, we obtain the pair criterion, firstly stated by {\v{C}ern\'y~\cite[Thm.~2]{Cer64}}.

\IN{pair criterion}
\begin{theorem}[Pair Criterion, {\v{C}ern\'y~\cite[Thm.~2]{Cer64}}]
	\label{prelim:thm:pair_sync_condition}
	An automaton $\semiautomat{A}{Q}{\Sigma}{\delta}$  with non-empty~$Q$
	is synchronizing if and only if for all pairs of states $p, q \in Q$
	there exists $u \in \Sigma^*$
	such that $\applytrans{\delta}{p}{u} = \applytrans{\delta}{q}{u}$. 
\end{theorem}
\begin{proof}
	By \cref{prelim:subclasses:sync:Aiffpairaut} 
	we have $\Syn(\mathcal A) \ne \emptyset$
	if and only if 
	$\Syn\left(\mathcal A^{\le 2} / {\sim_F}\right) \ne \emptyset$ where $F = \{\{q\} \mid q \in Q\}$.
	The pair automaton $\mathcal A^{\le 2} / {\sim_F}$
	has $F$ as a trap state.
	So, by \crefitem{prelim:thm:sync_trap_state}{prelim:subclasses:thm:trap:item:unique-max},
	we have
	$\Syn\left(\mathcal A^{\le 2} / {\sim_F}\right) \ne \emptyset$
	if and only if the trap state is reachable from every other state, i.e.,
	for $p, q\in Q$ there exists $u \in \Sigma^*$
	such that $\delta(\{p,q\}, u) \in F$, or equivalently
	$\delta(p, u) = \delta(q, u)$.	
\end{proof}

\begin{remark}
	In the language of transformation semigroup as discussed in \cref{prelim:remarks_trans_semigrp_semiautomata}, 
	the pair criterion states that a transformation semigroup $S$ on a finite set
	contains a constant map if and only if every two points can be mapped
	to a common point by some transformation from~$S$.
	The pair automaton has a trap state, but note that in general,
	we have a quadratic lower bound for trap-state automata, see \cref{prelim:fig:trap_state_tight}.	
\end{remark}

The following equivalences follow readily from \cref{prelim:thm:pair_sync_condition}.

\begin{corollary}
 Let $\semiautomat{A}{Q}{\Sigma}{\delta}$.
 Then the following are equivalent.
 \begin{enumerate}
 \item $\mathcal A$ is synchronizing.
 
 \item $\mathcal A \times \mathcal A$ is synchronizing.
 
 \item $\mathcal A^{\le 2}$ is synchronizing.
 
 \item $\mathcal A^{\le 2} / {\sim_F}$ with $F = \{\{q\} \mid q \in Q\}$ is synchronizing.
 
 \item $( \mathcal A \times \mathcal A ) / {\sim_D}$ with $D = \{(q,q) \mid q \in Q\}$ is synchronizing.
 \end{enumerate} 	
\end{corollary}

\IN{stability relation} \label{def:stability}
The pair criterion can be expressed in the following form using
the \emph{stability relation} $\sim_s$ introduced in~\cite{DBLP:journals/ijfcs/CulikKK02,DBLP:journals/jucs/Kari02}
where we set for $p,q \in Q$
\[
p \sim_s q \Leftrightarrow \forall u \in \Sigma^* \  \exists v \in \Sigma^* \colon \applytrans{\delta}{p}{uv} = \applytrans{\delta}{q}{uv},
\]
This is in fact a $\TS{\mathcal A}$-congruence.

\begin{corollary}[{Kari~\cite{DBLP:journals/jucs/Kari02,DBLP:journals/ijfcs/CulikKK02}}] \label{prelim:cor:stability-sync}
	An automaton $\semiautomat{A}{Q}{\Sigma}{\delta}$ with $Q \ne \emptyset$
	is synchronizing if and only if 
	the stability relation $\sim_s$ equals the universal relation.
\end{corollary}

Later, we will also use a transformation semigroup variant of this relation.
In this case, for a transformation semigroup $S$ on a domain~$X$
and $a,b \in X$ we set $a \sim_s b$
if and only if for all $t \in S$ there exists $h \in S$ such that $h(t(a)) = h(t(b))$.
Then~$S$ contains a constant transformation if and only if $\sim_s$ equals
the universal relation (see also \cref{prelim:remarks_trans_semigrp_semiautomata} for a discussion
of the relation of transformation semigroups and automata).

\subsection{Algorithmic Questions} 
\label{prelim:sync:algorithmic}

%!TEX root = ../../../thesis-trier.tex

First, in \cref{prelim:subclasses:sync:alg-unique-trap,prelim:subclasses:sync:thm:alg-unique-trap},
we investigate the running time of the problem of deciding synchronizability for automata with a trap state.
Then, \cref{thm:test_sync,prelim:thm:compute_sync_word}
make algorithmic statements for general automata. 
Lastly, in \cref{prelim:thm:comp_sync_states_if_sync},
we prove that for synchronizing automata
we can compute in linear time the set of synchronizing
states. As a consequence, stated in \cref{prelim:cor:set_sync_words_in_P},
for general automata the set of synchronizing states can be computed in quadratic time.

Later, in \cref{compl:sec:problems}, we will introduce a decision variant of the problem of computing a shortest synchronizing word.
In general, this problem is \NP-complete even for partially ordered binary automata, see~\cref{prelim:compl:sec:discussion}.

\begin{proposition}\label{prelim:subclasses:sync:alg-unique-trap}
	Let $\semiautomat{A}{Q}{\Sigma}{\delta}$.
	We can test in time $O(|Q|)$ and space $O(1)$
	whether $\mathcal A$ has a trap state, and also whether it has a unique trap state
	and compute this state if it exists.
\end{proposition}
\begin{proof}
	Iterating through all states and checking whether
	we find precisely one trap state can be done in the mentioned time and space bounds.
\end{proof} 

By \crefitem{prelim:thm:sync_trap_state}{prelim:subclasses:thm:trap:item:unique-max} 
an automaton with a trap state
is synchronizing precisely if this trap state is reachable
from every other state, which gives the following time and space bounds
to test synchronizability for automata with a trap state
that were also noted by Martyugin~\cite[Prop.~3.2]{Mar2009b}.

\begin{theorem}[{Martyugin~\cite[Prop.~3.2]{Mar2009b}}]\label{prelim:subclasses:sync:thm:alg-unique-trap}
	Let $\semiautomat{A}{Q}{\Sigma}{\delta}$ 
	be an automaton with a trap state.
	Then we can test in time $O(|Q| + |Q|\cdot |\Sigma|)$
	and space $O(|Q|)$ whether it is synchronizing.
\end{theorem} 
\begin{proof}
	By \cref{prelim:subclasses:sync:alg-unique-trap} 
	we can test whether the automaton has a unique trap, and if so compute it, in time $O(|Q|)$ and space $O(1)$.
	If this is the case, then checking whether there exists a path from every other state
	to this trap state can be done in time $O(|Q| + |\Sigma||Q|)$
	and space $O(|Q|)$ with a breadth-first search
	by following the edges in the automaton graph 
	in reverse order starting from the trap state (see, for example, the textbook~\cite{cormen01introduction} for an explanation of this algorithm
	and a statement of its resource requirements).	
\end{proof}

\begin{corollary} \label{prelim:subclasses:sync:cor:alg-unique-trap}
 For a class of automata with the property that a synchronizing state is a trap state,
 we can test whether an automaton from this class is synchronizing
 in time $O(|Q| + |Q|\cdot |\Sigma|)$
 and space $O(|Q|)$.
\end{corollary}
\begin{proof}
	First test in linear time and constant space whether an input automaton
	has a trap state, which can be done in the stated time and space bounds
	by \cref{prelim:subclasses:sync:alg-unique-trap}.
	If this is not the case, then it is not synchronizing by assumption
	Otherwise, by \cref{prelim:subclasses:sync:thm:alg-unique-trap},
	we can test in the given time and space bounds whether it is synchronizing.		
\end{proof}

By \cref{prelim:subclasses:sync:Aiffpairaut},
$\mathcal A$
is synchronizing if and only if
$\mathcal A^{\le 2} / {\sim_F}$ with $F = \{\{q\}\mid q \in Q\}$
is synchronizing.
As the pair automaton $\mathcal A^{\le 2} / {\sim_F}$ has the trap state~$F$,
$\frac{|Q|(|Q|-1)}{2} + 1$ states,
and $\left( \frac{|Q|(|Q|-1)}{2} + 1 \right) \cdot |\Sigma|$ transitions,
to test synchronizability of $\mathcal A$ we obtain the following time and space bounds from \cref{prelim:subclasses:sync:thm:alg-unique-trap}.

\begin{theorem}
	\label{thm:test_sync}
	Let $\semiautomat{A}{Q}{\Sigma}{\delta}$.
	We can test in  time $O\left(|Q|^2 + |Q|^2 \cdot |\Sigma|\right)$
	and space $O\left(|Q|^2\right)$	whether $\mathcal A$ is synchronizing.
\end{theorem}

However, the algorithm behind \cref{thm:test_sync} 
does not give us a synchronizing word. Natarajan~\cite[p.~138]{DBLP:conf/focs/Natarajan86}
has given a greedy algorithm to compute some synchronizing word of length at most $O(|Q|^3 \cdot |\Sigma|)$ in time $O(|Q|^4|\Sigma|)$.
Eppstein~\cite[Thm.~5, 6 \& 7]{Eppstein90} analyzed this algorithm further,
optimized it to run in $O(|Q|^3 + |Q|^2 \cdot |\Sigma|)$ time and $O(|Q|^2 + |Q| \cdot |\Sigma|)$ space, and proved that some synchronizing word of length at most $O(|Q|^3)$ 
can be computed by it (see also the surveys by Sandberg~\cite[Thm. 1.15]{Sandberg04}, Volkov~\cite[p.~16]{DBLP:conf/lata/Volkov08} 
and Kari \& Volkov~\cite[p.~531]{KariVolkov2021}).
Later, in \cref{compl:sec:subclasses:known-results}, we will also prove
that the deciding whether there exists a synchronizing word bounded by some given length-parameter
is \NP-complete.

\begin{theorem}[Eppstein~\cite{Eppstein90}, Natarajan~\cite{DBLP:conf/focs/Natarajan86}]
	\label{prelim:thm:compute_sync_word}
	Let $\semiautomat{A}{Q}{\Sigma}{\delta}$.
	We can test in time $O(|Q|^3 + |Q|^2 \cdot |\Sigma|)$
	and (not counting the space for the output)
	working space $O(|Q|^2 + |Q| \cdot |\Sigma|)$ whether $\mathcal A$
	is synchronizing and compute a synchronizing word of length $O(|Q|^3)$.
\end{theorem}

Now, for a synchronizing automaton, we can compute the set of synchronizing states in linear time. %time $O(|Q||\Sigma| + |Q|)$.

\begin{theorem}
	\label{prelim:thm:comp_sync_states_if_sync}
	Let $\semiautomat{A}{Q}{\Sigma}{\delta}$
	be a synchronizing automaton.
	Then we can compute in time $O(|Q| + |Q||\Sigma|)$ and space $O(|Q|)$
	the set of synchronizing states.	
\end{theorem}
\begin{proof}
	First, compute the strongly connected
	components, which can be done
	in the mentioned time and space bounds~\cite{cormen01introduction,DBLP:journals/siamcomp/Tarjan72}.
	Since~$\mathcal A$ is synchronizing,
	it follows from \cref{prelim:thm:SCC_sync_words} and \cref{prelim:prop:maximal-states-SCC}
	there exists
	a unique sink component, i.e., a strongly connected component
	with no outgoing transitions, 
	which can be easily identified and coincides with the set of synchronizing states.
	
	Alternatively, we can pick some state and start a breath-first search or a depth-first search
	in the automaton graph from this state.
	The first state $q$ we encounter with no unvisited 
	neighbors is a maximal state.
	Since~$\mathcal A$
	is synchronizing, $[q]_{\eqreach}$
	equals the set of synchronizing states
	by \cref{prelim:thm:SCC_sync_words}.
\end{proof}

For a synchronizing automaton with a trap state~$t$,
the set of synchronizing states equals $\{t\}$
and this set can be computed in linear time 
by \cref{prelim:subclasses:sync:alg-unique-trap}.

Combining Theorem~\ref{thm:test_sync} 
and~\ref{prelim:thm:comp_sync_states_if_sync}
gives the next result.

\begin{corollary}
	\label{prelim:cor:set_sync_words_in_P}
	Let $\semiautomat{A}{Q}{\Sigma}{\delta}$.
	We can compute in time 
	$
	O\left( |Q|^2 \cdot |\Sigma|\right)
	$
	and space $O\left(|Q|^2\right)$	
	the set of synchronizing states of $\mathcal A$.
\end{corollary}
\begin{proof}
	First, test whether $\mathcal A$ is synchronizing.
	If $\mathcal A$ is not synchronizing, then return the empty set.
	Otherwise, compute the set of synchronizing states and return
	this set. Combining the time and space bounds from Theorem~\ref{thm:test_sync} 
	and~\ref{prelim:thm:comp_sync_states_if_sync}
	for these computations gives the claim.
\end{proof}


